% GENERATED FILE. Edit canonical lessons, then run npm run generate:books.
% canonical-source-hash: fnv1a64:7115f7b0b2deb4bc
% canonical-lessons: TA-C10-vaara-kizhamai, TA-S05-rra, TA-W04-vowel-signs-nandri, TA-S114-vowel-sign-aa, TA-C10-checkpoint

\chapter{Days of the Week}
\label{ch:ta-days-of-the-week}
\hlchaptermodality{10}

\begin{chapteropening}
\textbf{By the end of this chapter:} \emph{I can name the seven weekdays, take each apart as planet plus \ta{கிழமை}, and say which are Tamil's own words and which arrived from Sanskrit.}

Complete TA-C10-checkpoint, the chapter's closing checkpoint: run the week from thiṅgaḷ to ñāyiṟu with kizhamai, sort the four native planet-words from the three Sanskrit ones, tell \ta{ந} from \ta{ன} and name \ta{ற}, find \ta{◌ா} in \ta{நான்}, then write \ta{ந}, \ta{ன}, \ta{ற} and \ta{நா} from sound.
\end{chapteropening}

\section[\texorpdfstring{\ta{திங்கள்}–\ta{ஞாயிறு} (thiṅgaḷ–ñāyiṟu)}{திங்கள்–ஞாயிறு (thiṅgaḷ–ñāyiṟu)}]{\ta{திங்கள்}–\ta{ஞாயிறு} — a home-grown and borrowed week}

\label{lesson:TA-C10-vaara-kizhamai}



\begin{quote}
Hindi's week is built entirely from Sanskrit deity-names plus \textbf{\dv{वार}} (\emph{vār}). Tamil's week tells a richer story: it uses its \textbf{own} word for \textquotedblleft{}day,\textquotedblright{} and \textbf{some} of its planet-names are Tamil's own too — while others are borrowed, just like Hindi's.
\end{quote}

\begin{grammarlens}[title={Grammar Lens: The pattern: {[}planet{]} + \ta{கிழமை}, \textquotedblleft{}day\textquotedblright{}}]
Every Tamil weekday ends in \textbf{\ta{கிழமை}} (\emph{kizhamai}), \textquotedblleft{}\textbf{day}\textquotedblright{} — a native Dravidian word, not the Sanskrit \textbf{\ta{வாரம்}} (\emph{vāram}) that Hindi, Kannada, and Telugu build on (though Tamil also keeps a more formal, Sanskritic \emph{-vāram} set alongside this everyday one):

\noindent
\begin{tabularx}{\linewidth}{@{}>{\raggedright\arraybackslash}X>{\raggedright\arraybackslash}X@{}}
\toprule
\textbf{Tamil weekday} & \textbf{planet-word, meaning, and source} \\
\midrule
\textbf{\ta{திங்கள்கிழமை}} & \emph{thiṅgaḷ}, \textquotedblleft{}\textbf{Moon}\textquotedblright{} — Tamil's \textbf{own} word \\
\textbf{\ta{செவ்வாய்க்கிழமை}} & \emph{cevvāy}, traditionally \textquotedblleft{}\textbf{the red one}\textquotedblright{} — likely Tamil's \textbf{own} descriptive name \\
\textbf{\ta{புதன்கிழமை}} & \emph{pudhaṉ}, \textbf{Mercury} — borrowed from Sanskrit \textbf{Budha} \\
\textbf{\ta{வியாழக்கிழமை}} & \emph{viyāzhaṉ}, \textbf{Jupiter} — likely from Sanskrit, via a name for \textbf{Bṛhaspati} \\
\textbf{\ta{வெள்ளிக்கிழமை}} & \emph{veḷḷi}, traditionally \textquotedblleft{}\textbf{silver}\textquotedblright{} — likely Tamil's \textbf{own} descriptive name \\
\textbf{\ta{சனிக்கிழமை}} & \emph{saṉi}, \textbf{Saturn} — borrowed from Sanskrit \textbf{Śani} \\
\textbf{\ta{ஞாயிற்றுக்கிழமை}} & \emph{ñāyiṟu}, \textquotedblleft{}\textbf{Sun}\textquotedblright{} — Tamil's \textbf{own} word \\
\bottomrule
\end{tabularx}
\end{grammarlens}

\begin{culture}
Notice the pattern: \textbf{Moon, Mars, Venus, and Sun} all carry Tamil's own descriptive or native words, while \textbf{Mercury, Jupiter, and Saturn} carry Sanskrit ones. This isn't Tamil \textquotedblleft{}failing\textquotedblright{} to have its own full week — it's a genuine \textbf{blend}, native Dravidian vocabulary standing right alongside Sanskrit loans, in the same list. (Note for the Tamil-speaking reviewer: this lesson spells Monday \textbf{\ta{திங்கள்கிழமை}}, the common form — the more traditional sandhi-shifted spelling \textbf{\ta{திங்கட்கிழமை}} may be preferred and should be swapped in on review if so. The \emph{cevvāy} \textquotedblleft{}red one\textquotedblright{} and \emph{viyāzhaṉ} Jupiter-name derivations in particular are worth treating as the traditional, commonly-given account rather than settled certainty — exact planet-name histories this old are not always fully nailed down.)
\end{culture}

\subsection*{Your turn}
Say these aloud:

\begin{itemize}
\raggedright
  \item \textquotedblleft{}thiṅgaḷ\textquotedblright{} — moon, Tamil's own word — then \textquotedblleft{}kizhamai,\textquotedblright{} day
  \item the borrowed ones — \textquotedblleft{}pudhaṉ, viyāzhaṉ, saṉi\textquotedblright{} — Mercury, Jupiter, Saturn from Sanskrit
  \item the native ones — \textquotedblleft{}thiṅgaḷ, cevvāy, veḷḷi, ñāyiṟu\textquotedblright{} — moon, Mars, Venus, sun
\end{itemize}

\begin{tcolorbox}[breakable,colback=teal!4,colframe=teal!35!black,title={Before you move on}]
What does every Tamil weekday end in, and where does that word come from? (\textbf{\ta{கிழமை}} \emph{kizhamai}, \textquotedblleft{}day\textquotedblright{} — Tamil's \textbf{own} word, not Sanskrit \emph{vāra}.) Which four planet-names are Tamil's own, and which three are Sanskrit borrowings? (\textbf{Moon, Mars, Venus, Sun are native}; \textbf{Mercury, Jupiter, Saturn are Sanskrit} — \emph{Budha, Bṛhaspati, Śani}.)
\end{tcolorbox}
\section[\texorpdfstring{\ta{ற} (ṟa)}{ற (ṟa)}]{\ta{ற} — the last of the family, and the one that drops below the line}

\label{lesson:TA-S05-rra}



\begin{quote}
Before the new one: say the sound of \textbf{\ta{ந}}.

One letter this time. Just one.

The fourth member of this top-bar family, and the one with a habit none of the others have: its right leg keeps going \textbf{below the baseline} and sweeps left into a long descender. Tamil letters mostly sit on the line. This one hangs off it.
\end{quote}

\subsection*{Script you'll notice: \ta{ற}}
\textbf{\ta{ற}} — \emph{ṟa}.

What it is made of:

\begin{itemize}
\raggedright
  \item two arches side by side, giving three legs down to the baseline
  \item the right leg continues BELOW the baseline, sweeps left, and drops into a long descender
\end{itemize}

It is the ṟ in \ta{நன்றி} — and with it, every consonant of \emph{thank you}.

\subsection*{Writing: \ta{ற}}
\hlblockfigure{\detokenize{figures/TA-S05-rra-filmstrip.pdf}}{How \ta{ற} is written, stroke by stroke}

\begin{itemize}
\raggedright
  \item \textbf{1.} start at the lower left, climb the left side, and arch to the middle
  \item \textbf{2.} without lifting, descend the middle upright
  \item \textbf{3.} without lifting, climb back up the same upright
  \item \textbf{4.} without lifting, arch over and descend the right side
  \item \textbf{5.} without lifting, sweep left below the baseline and drop into the long descender — and only now lift
\end{itemize}

\textbf{Pen lifts: 0.} The pen never leaves the paper.

\begin{quote}
Stroke order is one attested teaching order, not a national standard — Tamil handwriting is taught with school-to-school variation. Source: Sankaran Radhakrishnan, Tamil Script Learners Manual, Appendix I: Hand-movements, Frame 10, \ta{ற} (Univ. of Texas at Austin), p. 194.
\end{quote}

\subsection*{Your turn}
\begin{itemize}
\raggedright
  \item \emph{Trace it:} \ta{ற} once, slowly, following the steps above
  \item \emph{Say it:} \textquotedblleft{}ṟa\textquotedblright{} as you finish it
  \item \emph{Trace it:} it twice more, without looking back at the steps
\end{itemize}

\begin{tcolorbox}[breakable,colback=teal!4,colframe=teal!35!black,title={Before you move on}]
What does \ta{ற} do that \ta{ண}, \ta{ன} and \ta{ந} do not? (\textbf{Drops below the baseline.}) Which word have you now met all the consonants of? (\textbf{\ta{நன்றி}}.)
\end{tcolorbox}
\section[\texorpdfstring{\ta{ந}, \ta{ன}, \ta{ற} (na, ṉa, ṟa)}{ந, ன, ற (na, ṉa, ṟa)}]{\ta{ந}, \ta{ன}, \ta{ற} — the letter bodies inside \ta{நன்றி}}

\label{lesson:TA-W04-vowel-signs-nandri}



\begin{quote}
Lesson 2 showed you that Tamil has \textbf{three} n-letters and promised you'd draw the other two when a word needed them. Chapter 1's \emph{thank you} needs both.
\end{quote}

\subsection*{Script you'll notice: The two remaining n's}
\hlblockfigure{\detokenize{figures/TA-W04-vowel-signs-nandri-filmstrip.pdf}}{How these letters are written, one after another, stroke by stroke: \ta{ந}, \ta{ன}, \ta{ற}}

\textbf{\ta{ந}} — \emph{na}, the \textbf{dental} n, tongue against the teeth:

\begin{enumerate}
\raggedright
  \item \textbf{a straight top bar}
  \item \textbf{a vertical on the left}
  \item \textbf{a second vertical in the middle}
  \item \textbf{a right stroke that curls below the baseline} and sweeps left into a descender
\end{enumerate}

\textbf{\ta{ன}} — \emph{ṉa}, the \textbf{alveolar} n, tongue on the ridge behind the teeth:

\begin{enumerate}
\raggedright
  \item \textbf{a straight top bar}
  \item \textbf{a loop on the left}
  \item \textbf{one arch}
  \item \textbf{a straight vertical on the right, down to the baseline}
\end{enumerate}

Compare it with Lesson 2's \textbf{\ta{ண}}: same top bar, same loop, same closing vertical. \textbf{\ta{ண} simply has a second arch inserted before that vertical.} So \textbf{\ta{ன} is \ta{ண} minus one arch} — one arch for \ta{ன}, two for \ta{ண}. That is the whole contrast, and it is the cheapest way to keep them apart.

With Lesson 2's \textbf{\ta{ண}} (retroflex), that completes the set. Same consonant to an English ear; three letters, three tongue positions, in Tamil.

\subsection*{Script you'll notice: \ta{ற} — \textquotedblleft{}ṟa\textquotedblright{}, the hard r}
\begin{enumerate}
\raggedright
  \item \textbf{two arches side by side} — giving it three legs down to the baseline
  \item \textbf{the right leg keeps going below the baseline}, sweeps \textbf{left}, and drops into a long descender
\end{enumerate}

It is the tallest thing you have drawn: most Tamil letters sit on the baseline, and \ta{ற} hangs well beneath it.

The dot under \textbf{ṟ} marks it as the \emph{hard} or trilled r, distinct from Tamil's other r. You'll meet the contrast properly later; for now, learn the shape.

\subsection*{Your turn}
Write these out:

\begin{itemize}
\raggedright
  \item \ta{ந} — \textquotedblleft{}top bar, \textbf{two} verticals, then the curl that sweeps left\textquotedblright{}
  \item \ta{ன} — \textquotedblleft{}top bar, left loop, \textbf{one} arch, straight vertical\textquotedblright{}
  \item \ta{ற} — \textquotedblleft{}two arches, then the leg that sweeps left and drops\textquotedblright{}
\end{itemize}

\begin{tcolorbox}[breakable,colback=teal!4,colframe=teal!35!black,title={Before you move on}]
Name Tamil's three n-letters and their tongue positions. (\textbf{\ta{ந}} dental, \textbf{\ta{ன}} alveolar, \textbf{\ta{ண}} retroflex.) What are a consonant's three possibilities so far? (Keep the built-in \textbf{a} or \textbf{remove} it with the puḷḷi.) Which letter hangs below the baseline? (\textbf{\ta{ற}}, the hard \emph{ṟ}.) Next: the colours — black, white, red and blue.
\end{tcolorbox}
\section[\texorpdfstring{\ta{ா} (ā)}{ா (ā)}]{\ta{◌ா} — one character, met inside words you already say}

\label{lesson:TA-S114-vowel-sign-aa}



\begin{quote}
Before the new one: \ta{ல} — what does it do?

One character this time. Just one — and you have been saying it for pages without knowing which mark on the page it was.
\end{quote}

\subsection*{Script you'll notice: \ta{◌ா}}
\textbf{\ta{◌ா}} — \emph{ā}.

It is a \textbf{vowel sign}. It is not a letter and never stands alone: it attaches to a consonant and \textbf{replaces} the \emph{a} built into it with \emph{ā}. Replaces, not adds — the \emph{a} is gone.

Where it sits: a vertical stroke with a small top hook, written after the consonant.

Worked through: \textbf{\ta{ம}} + \textbf{\ta{◌ா}} = \textbf{\ta{மா}} — \emph{mā}.

You already say these, and every one of them has \ta{◌ா} somewhere inside it:

\begin{itemize}
\raggedright
  \item \textbf{\ta{வணக்கம்} / \ta{நமஸ்காரம்}} \emph{vaṇakkam / namaskāram} — Tamil kept a native bow-word where its neighbors borrowed Sanskrit
  \item \textbf{\ta{நான்}} \emph{nāṉ} — I
  \item \textbf{\ta{பரவாயில்லை}} \emph{paravāyillai} — it's okay / no problem / you're welcome
  \item \textbf{\ta{நாளை} \ta{பார்க்கலாம்}} \emph{nāḷai pārkkalām} — see you tomorrow
\end{itemize}

\subsection*{Writing: \ta{◌ா} — follow the numbered strip}
\hlblockfigure{\detokenize{figures/TA-S114-vowel-sign-aa-filmstrip.pdf}}{How \ta{ா} is written, stroke by stroke}

Follow the numbered strip: start where movement 1 begins, and let each panel show you the next movement before your pen makes it. Then write \ta{◌ா} yourself — slowly, and larger than it is printed.

\begin{quote}
The strip shows \textbf{where to start the character and which way to travel}, in one attested order. That is taught with real variation from school to school, so the strip names its source beneath it: treat it as a sound way in, not the only one.
\end{quote}

\subsection*{Your turn}
\begin{itemize}
\raggedright
  \item \emph{Look:} at these words, and find \ta{◌ா} in the ones that have it
\end{itemize}

\begin{quote}
\ta{நான்}  ·  \ta{வணக்கம்}
\end{quote}

\begin{itemize}
\raggedright
  \item \emph{Trace it:} \ta{◌ா} three times, saying \emph{ā} as you finish each one
  \item \emph{Look:} back at any page of this chapter and find \ta{◌ா} once more
\end{itemize}

\begin{tcolorbox}[breakable,colback=teal!4,colframe=teal!35!black,title={Before you move on}]
Which character is this — \ta{◌ா}? What vowel does it put on the consonant it attaches to — and what does it take off? (\textbf{Puts \emph{ā} on; takes the built-in \emph{a} off.}) Name one word you already say that contains it.
\end{tcolorbox}
\section[Practice]{Chapter 10 checkpoint — the week, three letters and a sign}

\label{lesson:TA-C10-checkpoint}



\begin{quote}
Nothing new here. The week out loud first, then the letters and the sign this chapter put on the page.
\end{quote}

\subsection*{Your turn — the week}
\begin{itemize}
\raggedright
  \item \emph{Say it:} the seven planet-words, Monday to Sunday, each followed by \emph{kizhamai}
\end{itemize}

(\emph{Thiṅgaḷ, cevvāy, pudhaṉ, viyāzhaṉ, veḷḷi, saṉi, ñāyiṟu}.)

\begin{itemize}
\raggedright
  \item What does \emph{kizhamai} mean, and whose word is it? (\textquotedblleft{}Day\textquotedblright{}: Tamil's own, not the Sanskrit \emph{vāram}.)
  \item Which four planet-words are Tamil's own, and which three came from Sanskrit? (Moon, Mars, Venus and Sun are native: \emph{thiṅgaḷ, cevvāy, veḷḷi, ñāyiṟu}. Mercury, Jupiter and Saturn are borrowed: \emph{pudhaṉ, viyāzhaṉ, saṉi}.)
\end{itemize}

\subsection*{Script — \ta{ந}, \ta{ன}, \ta{ற} and \ta{◌ா}}
\begin{itemize}
\raggedright
  \item Which n is \textbf{\ta{ந}}, and which is \textbf{\ta{ன}}? (\textbf{\ta{ந}} is the dental \emph{n}: top bar, two verticals, the curl that sweeps left. \textbf{\ta{ன}} is the alveolar \emph{ṉ}: \textbf{\ta{ண}} minus one arch.)
  \item Which letter hangs below the baseline, and whose consonants does it complete? (\textbf{\ta{ற}}, the hard \emph{ṟ}; with it you have every consonant of \textbf{\ta{நன்றி}}.)
  \item In \textbf{\ta{நான்}}, which mark turns the \emph{a} of \textbf{\ta{ந}} into \emph{ā}, and where does it sit? (\textbf{\ta{◌ா}}, written after the consonant.)
\end{itemize}

\emph{Read it:} \textbf{\ta{நான்}} aloud, and point to the \ta{◌ா} in it

\subsection*{Writing — from sound}
\emph{Cover:} the Script block above, so no model stays in view

\emph{Write it:} one letter each for the dental \emph{n}, the alveolar \emph{ṉ} and the hard \emph{ṟ}

\emph{Write it:} \emph{nā} — the dental \emph{n} carrying the \emph{ā} sign

\emph{Check:} against the key — \textbf{\ta{ந}}, \textbf{\ta{ன}}, \textbf{\ta{ற}}, \textbf{\ta{நா}}

\emph{Write it:} one repaired shape, then stop

\begin{tcolorbox}[breakable,colback=teal!4,colframe=teal!35!black,title={Before you move on}]
Say today's day in Tamil, then say whether its planet-word is Tamil's own or borrowed. That is the chapter: a week that blends two sources, and the letters that finish \emph{thank you}.

Next chapter: the colours.
\end{tcolorbox}
